\scanpage{13}
% archive sheet number: 12

(antiisomorphisme d'ens ordonnés).

\uline{Corollaire 3} Soit $\Phi$ un polyèdre combinatoire strict. Alors

a) $\Phi^{\unsure{\circ|o}}$ est un pol. combinatoire strict

b) Si $x,y\in\Phi$, avec $x\leq y$, alors l'ens. ordonné $\Phi_{xy}=\deleted{\unsure{f\ \colon\ x\leq z}}\{z\in\Phi\mid x\unsure{<|\leq}z\leq y\}$ est un polyèdre combinatoire strict.

\uline{Dém} a) Résulte du Cor. 2. On sait que $\Phi_{0,\unsure{x|y}}$ \margin{\begin{tikzpicture}[x=1em,y=1em]\draw[->] (0,.7) .. controls (-.1,0) and (.15,.12) .. (.55,.12);\end{tikzpicture} plus petit \unsure{élém|élém$^{\mathrm t}$|élt}} est un polyèdre combinatoire strict, donc par dualité aussi $\Phi_{x,1}$ \margin{\begin{tikzpicture}[x=1em,y=1em]\draw[->] (0,.35)--(.05,.05)--(.95,.08);\end{tikzpicture} pls grd él} 

$=(\Phi^{\circ}_{0,x})^\circ$. Appliquant ce dernier résultat à $x\in\Phi_{\unsure{0},y}$, on trouve que $\Phi_{xy}$ est un polyèdre combinatoire \unsure{stricte|strict}.

\deleted{\uline{\unsure{Cor 4|Cor. 4}} Soit \unsure{$x,y$} deux facettes}

\deleted{\uline{Cor. 4} \deleted{Soit $F$ une facette d'un p.} $C$ un polyèdre convexe. $F$ une facette \unsure{du dit|dudit}, $i$ \inserted{$\in\mathbb Z$} tel que $-1\leq i\leq d=\dim F$. Alors $\exists$ facette \deleted{de} $F'$ \unsure{incidente} à $F$ telle que $\dim F'=\unsure{i|.|{}}$}
